\section{Why width-thresholding works: covariance structure and step geometry}
\label{sec:mechanism}

The empirical ordering in Section~\ref{sec:results} poses two mechanistic
questions that loss curves alone cannot settle. First, why does diagonalizing
the width-2048 MLP-expansion output moments (\mlpdiag{}) match or slightly
exceed full spectral roots (\dense{}) while omitting those same moments
(\texttt{inner\_outer\_hidden}) forfeits half the gain and diagonalizing the
width-512 moments (\alldiag{}, \texttt{inner\_outer\_all\_diag}) degrades
quality? Second, does applying a unit-RMS outer factor $\hat L$ after the
finite polynomial map $\Nten$ genuinely rotate the update direction, or does it
merely inflate effective step norms at the fixed learning rate? Both questions
can be answered directly from the saved distributed checkpoints without
additional training.

We audited all 32 layers (eight ranks, 224 output moments and 128 input
moments, 352 symmetric FP32 EMA matrices in total) from the saved
\dense{} checkpoint at update 2{,}500 (mid-training), with matched
12-layer cross-checks at updates 500 and 5{,}120, alongside reconstructed
update directions across \muont{}, \mlpdiag{}, \alldiag{}, and \dense{}.
Table~\ref{tab:mechanism} and Figure~\ref{fig:mechanism} summarize the
measurements.

\begin{table}[t]
\centering
\caption{Covariance structure and diagonal preconditioner fidelity across all
32 layers at update 2{,}500 (\dense{} checkpoints, 352 FP32 EMA matrices).
``Diag.\ \%'' is $\|\diag(F)\|_2^2 / \|F\|_F^2$; $\mathrm{CV}(D)$ is the
coefficient of variation of the diagonal entries; $\lambda_{\max}/\bar\lambda$
and $\lambda_{\max}(R)$ are the top-to-mean eigenvalue ratios before and after
diagonal whitening ($R = D^{-1/2}FD^{-1/2}$). Relative errors compare
$L_{\diag}$ and $I$ against the exact spectral inner root $L = A_d(F)^{-1/2}$,
and $R^2(L) = 1 - \|L - L_{\diag}\|_F^2 / \|L - I\|_F^2$ is the fraction of
the preconditioner's departure from identity explained by diagonal scaling
alone.}
\label{tab:mechanism}
\scriptsize\setlength{\tabcolsep}{3.2pt}
% Generated by scripts/analyze_mechanism.py. Do not edit.
\begin{tabular}{lrrrrrrrrrr}
\toprule
Moment & $d$ & Diag.\ \% & $\mathrm{CV}(D)$ & $\lambda_{\max}/\bar\lambda$ & $\lambda_{\max}(R)$ & $\cos(L,L_{\diag})$ & $\mathrm{err}(L_{\diag})$ & $\mathrm{err}(I)$ & $R^2(L)$ \% & $R^2(\hat L)$ \% \\
\midrule
$\Out$: \texttt{gate} & 2048 & $81.9$ & $7.06$ & $386$ & $31.8$ & $0.9896$ & $0.113$ & $0.432$ & \textbf{$93.2$} & \textbf{$84.7$} \\
$\Out$: \texttt{up} & 2048 & $75.5$ & $5.07$ & $285$ & $48.5$ & $0.9826$ & $0.173$ & $0.452$ & \textbf{$85.4$} & \textbf{$75.1$} \\
\addlinespace
$\Out$: \texttt{q} & 512 & $22.6$ & $1.41$ & $68$ & $27.6$ & $0.9457$ & $0.369$ & $0.657$ & $66.4$ & $49.9$ \\
$\Out$: \texttt{k} & 512 & $17.4$ & $2.23$ & $107$ & $31.4$ & $0.9464$ & $0.357$ & $0.699$ & $72.1$ & $42.0$ \\
$\Out$: \texttt{v} & 512 & $24.4$ & $0.69$ & $80$ & $58.4$ & $0.8874$ & $0.580$ & $0.682$ & $26.5$ & $14.6$ \\
$\Out$: \texttt{o} & 512 & $41.7$ & $0.18$ & $42$ & $33.2$ & $0.9628$ & $0.313$ & $0.320$ & $3.0$ & $-3.7$ \\
$\Out$: \texttt{down} & 512 & $40.6$ & $0.17$ & $40$ & $32.1$ & $0.9665$ & $0.299$ & $0.305$ & $2.9$ & $-3.6$ \\
\addlinespace
$\In$: \texttt{down} & 2048 & $40.0$ & $5.44$ & $431$ & $69.9$ & $0.9766$ & $0.242$ & $0.515$ & $77.4$ & $55.3$ \\
$\In$: \texttt{attn} & 512 & $6.7$ & $1.29$ & $179$ & $113.0$ & $0.9188$ & $0.531$ & $0.571$ & $14.2$ & $-6.3$ \\
$\In$: \texttt{o} & 512 & $8.0$ & $0.84$ & $133$ & $119.7$ & $0.8851$ & $0.611$ & $0.725$ & $28.6$ & $27.2$ \\
$\In$: \texttt{mlp} & 512 & $4.7$ & $0.82$ & $168$ & $113.9$ & $0.9165$ & $0.531$ & $0.555$ & $8.8$ & $-9.5$ \\
\bottomrule
\end{tabular}

\end{table}

\subsection{Pre-SwiGLU Hadamard decorrelation versus residual rotational structure}
\label{sec:mech-hadamard}

Section~\ref{sec:freshness} noted that dense width-2048 eigenbases lag by up to
four updates, offering staleness as a candidate explanation for \mlpdiag{}'s
advantage over \dense{}. Table~\ref{tab:mechanism} exposes a cleaner,
architectural cause: \emph{pre-SwiGLU expansion cotangents are coordinate-aligned,
whereas residual-stream cotangents and RMSNormed inputs are rotation-dominated.}

Across all 32 layers, the width-2048 output moments $B_{\text{gate}}$ and
$B_{\text{up}}$ concentrate \mechOutGateDiagEnergyPct\% and
\mechOutUpDiagEnergyPct\% of their squared Frobenius norm on the diagonal
(medians \mechOutGateDiagEnergyMedPct\% and \mechOutUpDiagEnergyMedPct\%).
After normalizing into correlation matrices $R = D^{-1/2}BD^{-1/2}$, their
median off-diagonal root-mean-square correlation across the 32 layers is
\mechOutGateCorrRmsMed{} and \mechOutUpCorrRmsMed{} (bulk Layers 4--31 means:
\mechOutGateBulkCorrRms{} and \mechOutUpBulkCorrRms{}), sitting directly at the
single-update finite-sample floor $N_t^{-1/2} = 0.0052$ for $N_t = 36{,}608$
sampled tokens. Consequently, the elementwise diagonal root $L_{\diag}$ aligns
with the full spectral root $L$ at cosine \mechOutGateCosL{} and
\mechOutUpCosL{} (medians \mechOutGateCosLMed{} and \mechOutUpCosLMed{}) and
captures \mechOutGateExpVarLPct\% and \mechOutUpExpVarLPct\% of $\|L - I\|_F^2$
(bulk Layers 4--31: \mechOutGateBulkExpVarLPct\% and
\mechOutUpBulkExpVarLPct\%). Fitting $2.1\times 10^6$ off-diagonal entries per
expansion matrix therefore estimates finite-sample noise rather than signal;
zeroing those off-diagonals acts as exact shrinkage at zero basis age.

At the same time, $B_{\text{gate}}$ and $B_{\text{up}}$ are far from isotropic.
They exhibit the largest per-coordinate variance spread anywhere in the
transformer, with mean diagonal coefficients of variation
$\mathrm{CV}(D) = \mechOutGateDiagCv$ and $\mechOutUpDiagCv$ and max-to-mean
coordinate ratios of $\mechOutGateDiagMaxMean\times$ and
$\mechOutUpDiagMaxMean\times$. Omitting them in \texttt{inner\_outer\_hidden}
replaces $L$ by $I$, incurring relative error \mechOutGateDropErrL{} and
\mechOutUpDropErrL{} --- nearly four times larger than $L_{\diag}$'s
\mechOutGateRelErrL{} and \mechOutUpRelErrL{}.

\begin{finding}[Why width-thresholded preconditioning matches the architecture]
Pre-SwiGLU cotangents $\delta_{\text{up}} = (\delta_{\text{down}}W_{\text{down}})
\odot \operatorname{SiLU}(g)$ and $\delta_{\text{gate}} = (\delta_{\text{down}}
W_{\text{down}}) \odot (u \odot \operatorname{SiLU}'(g))$ are elementwise
Hadamard products of a zero-mean backward residual cotangent with a forward
gate signal. Under Schur multiplication of weakly coupled covariances
($\Sigma_{ij} \approx \Sigma^{\text{bwd}}_{ij}\Sigma^{\text{fwd}}_{ij}$), small
off-diagonal correlations multiply toward zero while per-neuron variances
multiply together, simultaneously decorrelating off-diagonals and amplifying
coordinate scale anisotropy. By contrast, the post-SwiGLU input moment
$C_{\text{down}}$ at the \emph{same} width 2048 has positive channel means from
one-sided $\operatorname{SiLU}$, leaving only \mechInDownDiagEnergyPct\%
diagonal energy and $54\times$ higher median correlation energy per coordinate
(\mechInDownCorrEnergyMed{} versus \mechOutGateCorrEnergyMed{}).
\end{finding}

The width-512 statistics invert this regime completely:
\begin{itemize}
\item \textbf{Residual-stream outputs (\texttt{o}, \texttt{down}) and value
  heads (\texttt{v}):} Because \texttt{o} and \texttt{down} project directly
  into the shared residual stream, their per-coordinate variances are nearly
  uniform ($\mathrm{CV}(D) = \mechOutODiagCv$ and $\mechOutDownDiagCv$).
  Diagonal scaling explains only \mechOutOExpVarLPct\% and
  \mechOutDownExpVarLPct\% of $\|L - I\|_F^2$ (\mechOutOExpVarLhatPct\% and
  \mechOutDownExpVarLhatPct\% on $\hat L$), missing $97\%$ of their rotational
  conditioning. Query and key cotangents (\texttt{q}, \texttt{k}) sit in
  between at \mechOutQExpVarLPct\% and \mechOutKExpVarLPct\%. This explains
  both why \alldiag{} retains two-thirds of \mlpdiag{}'s gain (via
  \texttt{gate}, \texttt{up}, \texttt{q}, \texttt{k}) and why \mlpdiag{}'s
  dense width-512 roots earn the remaining third (via \texttt{o},
  \texttt{down}, \texttt{v}).
\item \textbf{Width-512 inputs (\texttt{attn}, \texttt{o}, \texttt{mlp}):}
  Preceded by RMSNorm or multi-head attention, $92\%$--$95\%$ of $\|C\|_F^2$
  lies off the diagonal (Diag.\ \% of \mechInMlpDiagEnergyPct\% to
  \mechInODiagEnergyPct\%), and $R_{\diag}$ explains only
  \mechInMlpExpVarLPct\% to \mechInOExpVarLPct\% of $\|R - I\|_F^2$. Replacing
  input whitening with diagonals in \texttt{inner\_outer\_all\_diag} therefore
  destroys the primary conditioning of $\Nten$, producing the $+0.0180$-nat
  regression in Table~\ref{tab:main}.
\end{itemize}

\begin{figure}[t]
\centering
\begin{tikzpicture}
\begin{groupplot}[
  group style={group size=2 by 1, horizontal sep=1.6cm},
  width=0.46\textwidth, height=5.8cm,
  grid=major, grid style={gray!18},
  axis line style={gray!55}, tick align=outside,
  label style={font=\small}, tick label style={font=\scriptsize},
]
\nextgroupplot[
  ybar, bar width=7pt,
  ylabel={Explained variance $R^2(L)$ (\%)},
  ymin=-5, ymax=105,
  xtick={0,1,2,3,4,5,6,7,8,9,10},
  xticklabels={gate,up,q,k,v,o,down,in:dn,in:at,in:o,in:mlp},
  xticklabel style={rotate=40, anchor=east, font=\tiny},
  title={Diagonal fidelity across moments}, title style={font=\small},
]
\addplot[fill=laBlueLine!75, draw=laBlueLine!95]
  table[x=idx, y=exp_var_L_pct] {figs/mechanism_bars.dat};
\draw[gray!60, dashed] (axis cs:-0.6,0) -- (axis cs:10.6,0);

\nextgroupplot[
  unbounded coords=discard,
  xlabel={Projection ($m\times n$)},
  ylabel={Rotation angle vs.\ MuonT (deg)},
  ymin=0, ymax=31,
  xmin=-0.5, xmax=6.5,
  xtick={0,1,2,3,4,5,6},
  xticklabels={gate,up,q,k,v,o,down},
  xticklabel style={rotate=30, anchor=east, font=\scriptsize},
  title={\mlpdiag{} step rotation at norm $\mechMlpdiagGraftNorm{}$},
  title style={font=\small},
]
\addplot[ybar, bar width=10pt, fill=laPurpleLine!75, draw=laPurpleLine!95]
  table[x=idx, y=graft_deg_mlpdiag] {figs/mechanism_bars.dat};
\draw[laBlueLine, densely dashed, line width=0.9pt]
  (axis cs:-0.5,\mechMlpdiagGraftDeg) -- (axis cs:6.5,\mechMlpdiagGraftDeg);
\node[font=\tiny, text=laBlueLine!85!black, anchor=south east]
  at (axis cs:6.4,\mechMlpdiagGraftDeg) {mean $\mechMlpdiagGraftDeg^\circ$};
\end{groupplot}
\end{tikzpicture}
\caption{Checkpoint mechanism audit across all 32 layers at update 2{,}500.
\textbf{Left:} share of $\|L - I\|_F^2$ captured by diagonal scaling alone.
Width-2048 pre-SwiGLU cotangents (\texttt{gate}, \texttt{up}) sit at
\mechOutGateExpVarLPct\% and \mechOutUpExpVarLPct\%, whereas residual-stream
cotangents (\texttt{o}, \texttt{down}) sit at $3\%$ and width-512 inputs at
$9\%$--$29\%$. \textbf{Right:} angular rotation between \mlpdiag{}'s update
direction $D = \hat L\Nten(LUR)\hat R$ and the input-only direction
$D_{\text{in}} = \Nten(UR)\hat R$ on identical $(U,R,\hat R)$ states; the
simultaneous Frobenius norm ratio $\|D\|_F / \|D_{\text{in}}\|_F$ is
\mechMlpdiagGraftNorm{} ($\mechMlpdiagGraftNormPct\%$), ruling out step-size
inflation.}
\label{fig:mechanism}
\end{figure}

\subsection{Direction rotation rather than step-size inflation}
\label{sec:mech-graft}

Because $\Nten$ is not an exact polar factor, multiplying on the left by a
unit-RMS factor $\hat L$ does not algebraically preserve $\|D\|_F$. If
$\|D\|_F$ systematically exceeded $\|D_{\text{in}}\|_F = \|\Nten(UR)\hat R\|_F$,
output-side normalization at a fixed learning rate could act as an implicit
learning-rate increase. Evaluating both directions under exact BF16-operand
\texttt{tpu} arithmetic across all 224 matrices at updates 2{,}500 and 5{,}120
rules this out directly:
\begin{enumerate}
\item \textbf{Matched or slightly smaller direction norms:} Evaluated on the
  identical pre-update state $(U,R,\hat R)$, the instantaneous norm-grafting
  ratio $\|D\|_F / \|D_{\text{in}}\|_F$ is $\mechMlpdiagGraftNorm$
  ($\mechMlpdiagGraftNormPct\%$) for \mlpdiag{}, $\mechLeftdiagGraftNorm$
  ($\mechLeftdiagGraftNormPct\%$) for \alldiag{}, and $\mechDenseGraftNorm$
  ($\mechDenseGraftNormPct\%$) for \dense{} at update 2{,}500, and
  $\mechMlpdiagEndGraftNorm$ at update 5{,}120. Every projection kind lies
  within $[0.975, 0.996]$. Realized cross-trajectory direction norms against
  \muont{} are likewise $\mechMlpdiagRealDirNorm$ (\mlpdiag{}) and
  $\mechLeftdiagRealDirNorm$ (\alldiag{}).
\item \textbf{Nontrivial angular reorientation:} At that nearly identical norm,
  \mlpdiag{} rotates the update direction by a mean of
  $\mechMlpdiagGraftDeg^\circ$ ($\cos = \mechMlpdiagGraftCos$, spanning
  $17.0^\circ$--$26.1^\circ$ across all seven projections; Figure~\ref{fig:mechanism},
  right). In \alldiag{}, rotation on \texttt{o} and \texttt{down} collapses to
  $\mechLeftdiagOutODeg^\circ$ and $\mechLeftdiagOutDownDeg^\circ$ versus
  $\mechMlpdiagOutODeg^\circ$ and $\mechMlpdiagOutDownDeg^\circ$ in \mlpdiag{},
  matching Table~\ref{tab:mechanism}'s $3\%$ diagonal variance share on those
  two projections.
\item \textbf{Corroborating weight-norm trajectory:} Under fixed input-shaped
  weight decay $\lambda W\hat R^2$, an inflated effective learning rate would
  expand equilibrium weight norms $\|W\|_F$. Instead, total block weight norm
  at update 5{,}120 is $\mechMlpdiagWeightDiffEndPct\%$ \emph{smaller} than
  \muont{} in \mlpdiag{} and $\mechDenseWeightDiffEndPct\%$ smaller in
  \dense{}, while \nullmlp{} tracks \muont{} to $\mechNullMlpWeightDiffEndPct\%$.
\end{enumerate}
